%% cover_letter.tex
%% Cover letter for submission to Systems & Control Letters
%% NOT anonymised -- this is seen only by the editor, not reviewers

\documentclass[12pt]{article}
\usepackage[margin=2.5cm]{geometry}
\usepackage{amsmath,amssymb}
\usepackage{microtype}
\usepackage[colorlinks=true,urlcolor=blue,citecolor=blue]{hyperref}

\begin{document}

\noindent
Kunal Bhatia\\
Independent Researcher\\
Heidelberg, Germany\\
\href{mailto:kunal.bhatia@protonmail.com}{kunal.bhatia@protonmail.com}\\
ORCID: \href{https://orcid.org/0009-0007-4447-6325}{0009-0007-4447-6325}

\bigskip
\noindent\today

\bigskip
\noindent
Editor-in-Chief\\
\textit{Systems \& Control Letters}\\
Elsevier

\bigskip

\noindent Dear Editor-in-Chief,

\bigskip

I am pleased to submit the manuscript \textbf{``Task-Relative Descriptive
Privilege in Linear Gaussian Dynamical Systems''} for consideration in
\textit{Systems \& Control Letters}.

\paragraph{The result.}
The paper characterises, exactly and algebraically, when a single rank-$r$
linear observer can be simultaneously optimal for a family of squared-error
prediction targets in a finite-dimensional linear Gaussian dynamical system.
The central object is the \emph{target information matrix}
$M_{C,\tau} = \Sigma^{1/2}(A^\tau)^\top C^\top C A^\tau \Sigma^{1/2} \in \mathbb{R}^{n\times n}$,
whose top-$r$ eigenspace is the unique optimal observer for target $(C,\tau)$
(by the Ky Fan variational principle).
The main result is a biconditional: \emph{global privilege holds if and only
if all target information matrices share a common top-$r$ eigenspace} --- the
$r$-coherence condition.
Outside $r$-coherence, the optimal observer is task-relative, and the
worst-case regret gap $\Gamma > 0$ is computable from the closed-form Bayes
risk formula.

\paragraph{Scope and novelty.}
The Ky Fan principle and PCA optimality are classical.
The contribution is the exact algebraic characterisation of when a single
subspace suffices across a target family: $r$-coherence is an eigenspace
coincidence condition on the Grassmannian $\mathrm{Gr}(r,n)$, the natural
domain for rank-$r$ observer design.
The result is not asymptotic and does not rely on approximation.

\paragraph{Fit for the journal.}
The paper is a concise, self-contained theorem result in the linear systems
and observer design tradition of \textit{Systems \& Control Letters}.
The setting is the standard stationary linear Gaussian dynamical system
with squared-error prediction loss; the observer family is rank-$r$ linear
maps; the key structural quantity is the Grassmannian $\mathrm{Gr}(r,n)$.
The no-global-privilege result (Corollary 2) and the $r$-coherence equivalence
(Corollary 4) are directly relevant to observer selection and state estimation
under multiple objectives.
MSC 2020 codes: 93B07 (Observability); 93E10 (Estimation and detection);
15A18 (Eigenvalues, singular values, and eigenvectors).

\paragraph{Numerical illustration.}
Two illustrative instances are provided at $n=8$, $r=2$: a
no-global-privilege case (minimax gap $\Gamma = 0.026 > 0$, confirmed by
300-restart projected L-BFGS-B on $\mathrm{Gr}(2,8)$) and a structurally
coherent case (maximum pairwise principal angle $0.000^\circ$ across four
prediction horizons, robust to $\varepsilon=0.1$ perturbations).
All numbers are computed from the closed-form Bayes risk formula; there is
no Monte Carlo simulation.

\paragraph{Originality and prior publication.}
This manuscript has not been published elsewhere and is not currently under
consideration by any other journal.
Reproducible code (Python, NumPy/SciPy, $<$250 lines) is available at
\url{https://github.com/kunalb541/LGDS}.

\paragraph{Conflicts of interest.}
The author declares no conflicts of interest.

\bigskip

\noindent Thank you for considering this submission.

\bigskip

\noindent Sincerely,

\bigskip\bigskip

\noindent Kunal Bhatia

\end{document}
